\hypertarget{chapter-25-c17-standard-library-additions}{%
\section{CHAPTER 25: C17 STANDARD LIBRARY
ADDITIONS}\label{chapter-25-c17-standard-library-additions}}

\hypertarget{c17-standard-library-additions}{%
\section{C++17 STANDARD LIBRARY
ADDITIONS}\label{c17-standard-library-additions}}

\hypertarget{stdbyte}{%
\subsection{\texorpdfstring{1.
\texttt{std::byte}}{1. std::byte}}\label{stdbyte}}

A distinct type for byte-oriented memory access. Unlike
\passthrough{\lstinline!char!}, it is not an arithmetic type (prevents
accidental math).

\begin{lstlisting}[language={C++}]
#include <cstddef>

std::byte b = std::byte{0xAB};
// b += 1; // Error
int i = std::to_integer<int>(b);
\end{lstlisting}

\hypertarget{algorithms}{%
\subsection{2. Algorithms}\label{algorithms}}

\begin{itemize}
\tightlist
\item
  \passthrough{\lstinline!std::sample!}: Selects n random elements from
  a range.
\item
  \passthrough{\lstinline!std::clamp!}: Clamps a value between lo and
  hi.
\item
  \passthrough{\lstinline!std::gcd!},
  \passthrough{\lstinline!std::lcm!}: Math utilities.
\end{itemize}

\begin{lstlisting}[language={C++}]
int x = std::clamp(10, 0, 5); // 5
int g = std::gcd(12, 18);     // 6
\end{lstlisting}

\hypertarget{mathematical-special-functions}{%
\subsection{3. Mathematical Special
Functions}\label{mathematical-special-functions}}

Support for Laguerre polynomials, Bessel functions, elliptic integrals,
etc., in \passthrough{\lstinline!<cmath>!}.

\hypertarget{elementary-string-conversions-charconv}{%
\subsection{\texorpdfstring{4. Elementary String Conversions
(\texttt{\textless{}charconv\textgreater{}})}{4. Elementary String Conversions (\textless charconv\textgreater)}}\label{elementary-string-conversions-charconv}}

Low-level, allocation-free, locale-independent string conversions.
Extremely fast.

```cpp \#include

char buffer{[}10{]}; int value = 42;

// To chars std::to\_chars(buffer, buffer + 10, value);

// From chars std::from\_chars(buffer, buffer + 10, value); \# VOLUME
04: GODHOOD SUMMARY

\hypertarget{c17-landmark-features-reference}{%
\subsubsection{C++17 LANDMARK FEATURES
REFERENCE}\label{c17-landmark-features-reference}}

\begin{longtable}[]{@{}
  >{\raggedright\arraybackslash}p{(\columnwidth - 6\tabcolsep) * \real{0.25}}
  >{\raggedright\arraybackslash}p{(\columnwidth - 6\tabcolsep) * \real{0.25}}
  >{\raggedright\arraybackslash}p{(\columnwidth - 6\tabcolsep) * \real{0.25}}
  >{\raggedright\arraybackslash}p{(\columnwidth - 6\tabcolsep) * \real{0.25}}@{}}
\toprule
\# & Feature & Explanation & Code Example \\
\midrule
\endhead
1 & \textbf{Structured Bindings} & Unpack tuples, pairs, and structs
into named variables &
\passthrough{\lstinline!auto [x, y] = my\_pair;!} \\
2 & \textbf{if constexpr} & Compile-time conditional branching in
templates & \passthrough{\lstinline!if constexpr (is\_int\_v<T>)!} \\
3 & \textbf{Init-statements} & \passthrough{\lstinline!if!} and
\passthrough{\lstinline!switch!} can now initialize local variables &
\passthrough{\lstinline"if (auto it = m.find(k); it != m.end())"} \\
4 & \textbf{Fold Expressions} & Simplify variadic template unpacking
with operators & \passthrough{\lstinline!(... + args)!} \\
5 & \textbf{CTAD} & Class Template Argument Deduction from constructors
& \passthrough{\lstinline!std::vector v = \{1, 2, 3\};!} \\
6 & \textbf{Inline Variables} & Define static data members in headers
without ODR issues &
\passthrough{\lstinline!static inline int val = 5;!} \\
7 & \textbf{std::string\_view} & Non-owning, zero-allocation string
reference & \passthrough{\lstinline!void f(std::string\_view sv);!} \\
8 & \textbf{std::optional} & Type-safe representation of a maybe-present
value & \passthrough{\lstinline!std::optional<int> res;!} \\
9 & \textbf{std::variant} & Type-safe union (discriminated union) &
\passthrough{\lstinline!std::variant<int, float> v;!} \\
10 & \textbf{std::any} & Type-safe container for any single value &
\passthrough{\lstinline!std::any a = 42;!} \\
11 & \textbf{std::filesystem} & Standard library for file and directory
manipulation & \passthrough{\lstinline!fs::exists("path");!} \\
12 & \textbf{Parallel Algos} & Standard algorithms with execution
policies &
\passthrough{\lstinline!std::sort(std::execution::par, v.begin(), v.end());!} \\
13 & \textbf{std::scoped\_lock} & Deadlock-avoiding multi-mutex RAII
lock & \passthrough{\lstinline!std::scoped\_lock lk(m1, m2);!} \\
14 & \textbf{Guaranteed Copy Elision} & Compiler MUST omit copies in
specific return scenarios &
\passthrough{\lstinline!T f() \{ return T(); \}!} \\
15 & \textbf{std::byte} & Distinct type for raw memory bits &
\passthrough{\lstinline!std::byte b\{0xFF\};!} \\
\bottomrule
\end{longtable}

C++17 was the release of \textbf{Simplification and Vocabulary}. It
focused on making the language cleaner and providing standard types for
common patterns. 1. \textbf{Structured Bindings}: Unpacking tuples and
structs became trivial. 2. \textbf{if constexpr}: Compile-time branching
simplified template metaprogramming. 3. \textbf{Vocabulary Types}:
\passthrough{\lstinline!std::optional!},
\passthrough{\lstinline!std::variant!}, and
\passthrough{\lstinline!std::any!} replaced unsafe C-style patterns. 4.
\textbf{Filesystem}: Finally, a standard way to talk to the OS about
files.

\textbf{The Golden Rule of C++17}: Use
\passthrough{\lstinline!std::optional!} instead of null pointers, and
\passthrough{\lstinline!string\_view!} for efficient string passing. You
have simplified the vocabulary of your code.

\hypertarget{volume-05-gigantic-leap-c20}{%
\section{VOLUME 05 GIGANTIC LEAP
C20}\label{volume-05-gigantic-leap-c20}}

C++20 is the most significant update to the language since C++11. It
introduces the \textbf{Four Great Pillars} that fundamentally change how
we architect C++ software.

\hypertarget{the-four-great-pillars-head-first-style}{%
\subsubsection{The Four Great Pillars (Head First
Style)}\label{the-four-great-pillars-head-first-style}}

\begin{longtable}[]{@{}
  >{\raggedright\arraybackslash}p{(\columnwidth - 4\tabcolsep) * \real{0.33}}
  >{\raggedright\arraybackslash}p{(\columnwidth - 4\tabcolsep) * \real{0.33}}
  >{\raggedright\arraybackslash}p{(\columnwidth - 4\tabcolsep) * \real{0.33}}@{}}
\toprule
Pillar & Analogy & Why we need it \\
\midrule
\endhead
\textbf{Concepts} & \textbf{The Bouncer at the Club} & Before C++20,
templates were ``all are welcome.'' If you brought the wrong type, the
compiler would wait until you were inside the club to scream at you.
Concepts are like a bouncer at the door who checks your ID (type) before
you even enter. \\
\textbf{Modules} & \textbf{Sealed Folders vs.~Messy Desks} &
\passthrough{\lstinline!\#include!} is like dumping a giant pile of
messy blueprints on your desk every time you want to build a small part.
Modules are like sealed folders; you just grab exactly what you need
without making a mess of your current workspace. \\
\textbf{Coroutines} & \textbf{The Expert Chef} & A normal function is
like a chef who \emph{must} finish a whole recipe before doing anything
else. A Coroutine is a chef who can pause a recipe to wait for the oven
to heat up, work on another dish, and then come back exactly where they
left off. \\
\textbf{Ranges} & \textbf{The LEGO Pipe Factory} & Instead of manually
moving items from one box to another using iterators, Ranges let you
snap together ``pipes'' (filters, transforms) to create a high-speed
data assembly line. \\
\bottomrule
\end{longtable}

\begin{center}\rule{0.5\linewidth}{0.5pt}\end{center}
